\documentclass[11pt]{article}

\usepackage[margin=1in]{geometry}
\usepackage{booktabs}
\usepackage{graphicx}
\usepackage{amsmath}
\usepackage{amssymb}
\usepackage{hyperref}

\title{Learning Generative Strategy Spaces with Energy-Based Models and Game-Theoretic Evaluation for Inspectable Multi-Agent Reasoning}
\author{Research Scaffold Contributors}
\date{}

\begin{document}
\maketitle

\begin{abstract}
We present an early research system for separating high-level strategy generation from low-level action execution in adversarial multi-agent environments.
The system combines an energy-based model over latent strategy embeddings, approximate sampled-response evaluation, a strategy-conditioned policy, and lightweight update hooks for policy learning, contrastive energy learning, and one-step world-model fitting.
The current contribution is methodological and diagnostic: the framework makes strategy generation, evaluation, execution, and update behavior inspectable as separate components.
We explicitly report approximate metrics and negative diagnostic results; the present system does not claim state-of-the-art performance, Nash convergence, or exact exploitability.
\end{abstract}

\section{Introduction}

Multi-agent reinforcement learning systems often entangle strategic intent and low-level control in a single policy.
This makes it difficult to inspect whether failure comes from poor strategy generation, weak opponent evaluation, brittle execution, or unstable learning updates.
We study a decomposed alternative: generate candidate latent strategies, evaluate them against sampled opponent responses, execute the selected strategy with a conditioned policy, and update the components from rollout data.

The guiding research question is whether high-level strategies can be generated, evaluated, selected, and executed as a distinct layer above low-level policies.
This paper frames the contribution as an inspectable methods scaffold rather than a performance benchmark.

\section{Related Work}

The project sits at the intersection of energy-based generative models, population and best-response methods for games, latent option or skill discovery, and model-based reinforcement learning.
The related-work section should cover energy-based models for structured generation, PSRO-style population training, approximate exploitability and regret diagnostics, options or strategy embeddings, and world-model evaluation.

\section{Method}

\subsection{Generate, Evaluate, Execute, Update}

At each iteration, the system executes the following loop.
\begin{enumerate}
  \item \textbf{Generate}: sample latent strategy embeddings from an energy model using Langevin dynamics, optionally mixed with named heuristic strategies for debugging.
  \item \textbf{Evaluate}: score candidates against sampled opponent heuristics using average-case value, worst-case value, a robustness score, and an exploitability proxy.
  \item \textbf{Execute}: condition a policy on the selected strategy and roll it out in the attacker-defender gridworld.
  \item \textbf{Update}: store rollout summaries in a strategy buffer and apply lightweight policy, energy-model, and world-model updates.
\end{enumerate}

\subsection{Energy-Based Strategy Generator}

Let $z \in \mathbb{R}^d$ denote a latent strategy embedding and $E_\theta(z)$ denote a scalar energy.
Candidate strategies are sampled by unadjusted Langevin dynamics:
\[
z_{t+1} = z_t - \frac{\eta}{2}\nabla_z E_\theta(z_t) + \epsilon_t,\quad \epsilon_t \sim \mathcal{N}(0, \eta I).
\]
The current implementation uses contrastive divergence-style training with high-scoring buffer strategies as positives and sampled strategies as negatives.

\subsection{Sampled-Response Evaluation}

The evaluator is intentionally approximate.
For a candidate strategy, it rolls out against a finite set of heuristic opponent responses and records average-case value, worst-case value, robustness, and an exploitability proxy.
The exploitability proxy is a diagnostic over sampled responses only and must not be interpreted as Nash exploitability.

\subsection{Strategy-Conditioned Execution}

The policy receives both environment state and a latent strategy embedding.
This design lets the same low-level policy execute different high-level strategies while preserving a separate inspection surface for strategy selection.

\section{Experimental Setup}

\subsection{Primary Environment}

The main environment is a configurable attacker-defender gridworld.
The attacker starts near one corner and must reach a goal before a greedy defender intercepts it.
The environment is deliberately small so failures in generation, evaluation, and execution can be inspected directly.

\subsection{Baselines}

We compare the strategy loop against a uniform random attacker, a direct-to-goal heuristic, and a PPO-lite actor-critic attacker.
The direct-to-goal heuristic is expected to be strong in the default grid layout; its role is to prevent overclaiming from a weak learning baseline.

\subsection{Ablations}

The paper-facing ablation suite covers EBM versus Gaussian sampling, buffer positives versus random positives, robustness-aware versus average-value selection, candidate count, sampled opponent count, and world-model fitting on or off.

\section{Results}

Results should be reported as diagnostics.
The current public snapshot shows that the direct-to-goal heuristic wins the default gridworld while the early strategy loop and PPO-lite baseline do not.
This is evidence that the harness can expose failure modes, not evidence of robust multi-agent reasoning.

\begin{table}[h]
\centering
\caption{Current public baseline snapshot. Values come from the May 15 diagnostic run and should be replaced by multi-seed results before submission.}
\begin{tabular}{lrr}
\toprule
Baseline & Win rate & Episode return \\
\midrule
Direct-goal heuristic & 1.000 & 0.910 \\
PPO-lite baseline & 0.000 & -1.180 \\
Day-2 strategy loop & 0.000 & -1.220 \\
Random policy & 0.000 & -1.208 \\
\bottomrule
\end{tabular}
\end{table}

The fuller paper-facing results prose is maintained in \texttt{paper/results\_and\_limitations.tex} while the experiment suite is still evolving.

\section{Limitations}

The present system does not claim state-of-the-art performance, convergence to equilibrium, exact Nash exploitability, or robust multi-agent reasoning.
The evaluator uses sampled opponent responses, the world model is fit only on one-step transitions, and the pursuit/evasion infrastructure is currently a supplementary diagnostic track rather than an EBM-driven domain.

\section{Reproducibility}

The repository includes YAML experiment configs, unit and smoke tests, public JSON artifacts, figure scripts, an ablation runner, and a multi-seed protocol runner.
Before submission, paper tables should be regenerated from the scripted outputs rather than copied from hand-written snapshots.

\bibliographystyle{plain}
\bibliography{references}

\end{document}
